\documentclass[11pt]{article}
\usepackage[a4paper,margin=1in]{geometry}
\usepackage{amsmath,amssymb,amsthm}
\usepackage[hidelinks]{hyperref}
\usepackage{microtype}

\newtheorem{proposition}{Proposition}
\newcommand{\R}{\mathbb R}
\newcommand{\one}{\mathbf 1}

\title{A negative answer to the scaled-power\\order-continuity question}
\author{Literature-implied answer for arXiv:1208.3498}
\date{13 August 2026}

\begin{document}
\maketitle

\begin{center}
\fbox{\parbox{0.9\textwidth}{\raggedright\textbf{Status.}
\texttt{literature\_implied\_answer (negative answer)}.  Hadwin--Kitover--
Orhon's construction directly answers the question and also yields a
band-irreducible version.  This is an identification of known ingredients,
not a new solution claim.}}
\end{center}

\section{Source question}

For a positive operator $T$, Gao and Troitsky define $\R_+T$ to be the
smallest norm-closed semigroup containing $T$ and closed under multiplication
by nonnegative scalars.  Thus it contains every positive scalar multiple of
a power of $T$ and every norm limit
\[
             \lim_j b_jT^{n_j}\qquad(b_j\geq0,\ n_j\nearrow\infty).
\]
In Section~8 of \cite{GT}, PDF page 23, they ask whether
\[
 T\text{ positive and order-continuous}
 \quad\Longrightarrow\quad
 \R_+T\text{ consists of order-continuous operators}.
\]

\section{Direct literature implication}

Example~11 of Hadwin, Kitover, and Orhon \cite{HKO} constructs a Banach
lattice $X$ and a positive order-continuous operator $A$ on $X$.  Their
Remark~15 (supporting PDF, p.~15) records that
\[
                 A^n\longrightarrow P
                 \quad\text{in operator norm},
\]
where $P$ is a positive spectral projection that is not order-continuous
(not even $\sigma$-order-continuous).  Taking $b_j=1$ and $n_j=j$ in the
definition above gives $P\in\R_+A$.  Hence the answer to the source question
is negative.

The supporting paper does not explicitly identify the later wording in
\cite{GT}; the implication is direct but agent-identified.  This is why the
result is classified as a literature-implied answer.

\section{The band-irreducible context also fails}

The source question occurs inside a section on band-irreducible semigroups.
The same Example~11 yields the following stronger statement.

\begin{proposition}\label{prop:strong}
There are a Banach lattice $X$ and a positive, order-continuous,
band-irreducible operator $B$ such that $B^n$ converges in operator norm to
a positive projection that is not order-continuous.  In particular,
$\R_+B$ does not consist of order-continuous operators.
\end{proposition}

\begin{proof}
We use the notation and properties proved in Example~11 of \cite{HKO}.
The construction gives a Banach lattice
\[
 X=J\oplus\operatorname{span}\{\one\},\qquad
 J=\{x\in X:x(0)=0\},
\]
and a positive order-continuous lattice homomorphism $A$ satisfying
$A\one=\one$, $A(J)\subseteq J$, and
\begin{equation}\label{eq:factorial}
        \|(A|_J)^k\|\leq\frac1{k!}.
\end{equation}
It also constructs a positive order-continuous band-irreducible operator
$T$ and lets $L$ be multiplication by the coordinate function $t$.  The
paper proves that, for every sufficiently small $\varepsilon>0$,
\[
              B_\varepsilon=A+\varepsilon LT
\]
is positive, order-continuous, and band irreducible.

Evaluation at zero is fixed by $B_\varepsilon$: indeed
$(Ax)(0)=x(0)$ and $(LTx)(0)=0$.  Consequently $J$ is invariant and the
operator induced on $X/J$ is the identity.  Put
\[
 C_\varepsilon=B_\varepsilon|_J
 =A|_J+\varepsilon(LT)|_J.
\]
By \eqref{eq:factorial}, $r(A|_J)=0$.  Upper semicontinuity of the spectral
radius therefore permits us to choose $\varepsilon>0$ so small that
$r(C_\varepsilon)<1$, while retaining all the preceding properties.

Relative to $X=J\oplus\operatorname{span}\{\one\}$, write
\[
 B_\varepsilon=
 \begin{pmatrix} C_\varepsilon&d\\0&1\end{pmatrix}
 \quad(d\in J).
\]
Since $r(C_\varepsilon)<1$, we have $C_\varepsilon^n\to0$ and
$\sum_{k=0}^{n-1}C_\varepsilon^kd\to(I-C_\varepsilon)^{-1}d$.  Hence
\[
 B_\varepsilon^n\longrightarrow
 P=\begin{pmatrix}0&(I-C_\varepsilon)^{-1}d\\0&1\end{pmatrix}
 =u\otimes\delta_0,
\]
where $\delta_0(x)=x(0)$ and $u=P\one$ satisfies $u(0)=1$.

Remark~15 of \cite{HKO} notes that $J=\ker\delta_0$ is an ideal but not a
band; equivalently, $\delta_0$ is not order-continuous.  A nonzero positive
rank-one operator $u\otimes\delta_0$ is order-continuous only if
$\delta_0$ is order-continuous.  Therefore $P$ is not order-continuous,
although every $B_\varepsilon^n$ is.  Since $P\in\R_+B_\varepsilon$, the
stronger assertion follows.
\end{proof}

\section{Scope}

The proposition settles the precise closure implication negatively, even in
the band-irreducible setting.  It does not assert that all of the
Perron--Frobenius conclusions sought in \cite{GT} fail under every weaker
hypothesis; it only removes order continuity of the entire closed semigroup
as an automatic consequence of order continuity of its generator.

\begin{thebibliography}{9}
\bibitem{GT}
N. Gao and V. G. Troitsky,
\emph{Irreducible semigroups of positive operators on Banach lattices},
Linear Multilinear Algebra \textbf{62} (2014), 74--95;
arXiv:1208.3498.

\bibitem{HKO}
D. W. Hadwin, A. K. Kitover, and M. Orhon,
\emph{Strong monotonicity of spectral radius of positive operators},
Houston J. Math. \textbf{41}, no.~2;
arXiv:1205.5583v2.
\end{thebibliography}

\end{document}
